% ============================================================================
% Section 8.6 — Circle Graphs (Additional: FL, MN, TX, VA)
% ============================================================================
\section{Circle Graphs}

\begin{stepGoal}
\begin{itemize}
    \item Read and interpret circle (pie) graphs
    \item Convert between percents, fractions, and sector angles
\end{itemize}
\end{stepGoal}

\begin{stepVocabBox}
    \vocabItem{Circle Graph}{A graph where a circle is divided into \textbf{sectors} showing parts of a whole.}
    \vocabItem{Sector}{A ``pie slice'' — its size shows what \textbf{fraction} of the whole it represents.}
\end{stepVocabBox}

\begin{stepsCard}[How to Read and Use a Circle Graph]
    \stepItem{Find the \textbf{percent} (or fraction) each sector represents.}
    \stepItem{To find a \textbf{sector angle}: multiply the percent (as a decimal) by $360°$.}
    \stepItem{To find \textbf{how many}: multiply the percent (as a decimal) by the \textbf{total}.}
    \stepItem{Check: all percents must add to $100\%$ and all angles must add to $360°$.}
\end{stepsCard}

% --- Worked Example 1 ---
\begin{stepExample}{$200$ students were surveyed about their favorite subject: Math $30\%$, Science $25\%$, English $20\%$, Art $15\%$, Other $10\%$. How many chose Math? What angle for Science?}
    \stepShow{1} Math $= 30\%$. Science $= 25\%$.

    \stepShow{2} Science angle $= 0.25 \times 360° = 90°$.

    \stepShow{3} Students who chose Math $= 0.30 \times 200 = 60$.

    \stepShow{4} Check: $30 + 25 + 20 + 15 + 10 = 100\%$. \checkmark
    \stepResult{$60$ students chose Math; the Science sector is $90°$.}
\end{stepExample}

% --- Worked Example 2 ---
\begin{stepExample}{A sector in a circle graph has an angle of $72°$. What percent of the whole does it represent?}
    \stepShow{1} The sector angle is $72°$ out of $360°$ total.

    \stepShow{2} As a fraction: $\dfrac{72}{360} = \dfrac{1}{5}$.

    \stepShow{3} As a percent: $\dfrac{1}{5} = 0.20 = 20\%$.
    \stepResult{The sector represents $20\%$ of the whole.}
\end{stepExample}

\mascotStepTip{Remember: a full circle is $360°$. A quarter is $90°$, a half is $180°$.}

% ============================================================================
% PRACTICE
% ============================================================================
\newpage
\begin{stepPractice}{Circle Graphs Practice}
    \resetProblems

    \stepPracticeHeader[step1Color]{Find the Sector Angle}
    \begin{multicols}{2}
    \prob A category is $40\%$. Its angle? \answerBlank[2cm]
    \answer{$0.40 \times 360° = 144°$}
    \prob A category is $\frac{1}{8}$. Its angle? \answerBlank[2cm]
    \answer{$\frac{1}{8} \times 360° = 45°$}
    \end{multicols}

    \stepPracticeHeader[step2Color]{Find How Many}
    \prob $35\%$ of $300$ people prefer dogs. How many is that? \answerBlank[2.5cm]
    \answer{$0.35 \times 300 = 105$ people}

    \stepPracticeHeader[step3Color]{Find the Percent}
    \prob A sector has an angle of $54°$. What percent is it? \answerBlank[2.5cm]
    \answer{$\frac{54}{360} = 0.15 = 15\%$}

    \stepPracticeHeader[step4Color]{Check the Total}
    \wordProblem{A circle graph shows: Red $25\%$, Blue $30\%$, Green $20\%$, Yellow $x\%$. Find $x$ and the angle for Red.}{~}
    \answer{$x = 100 - 25 - 30 - 20 = 25\%$. Red angle $= 0.25 \times 360° = 90°$.}
\end{stepPractice}
